% !TeX root = ../main.tex
\addcontentsline{toc}{section}{Introduction}
\section*{Introduction}
Ce chapitre présente la conception de l'architecture agentique du diagnostic technique CCU. Il décrit le pipeline global orchestré par LangGraph, le rôle de chaque agent, le modèle de données partagé, le graphe Neo4j et le mécanisme de recherche augmentée par graphe (GraphRAG).

\section{Vue d'ensemble du pipeline}
L'architecture repose sur un graphe d'état (state graph) implémenté avec \textbf{LangGraph}. Chaque nœud du graphe correspond à un agent spécialisé qui lit et écrit dans un état partagé. Le flux de traitement d'un incident suit le cheminement suivant :

\begin{center}
\texttt{START $\rightarrow$ intake $\rightarrow$ collectors $\rightarrow$ root\_cause $\rightarrow$ ticket\_manager}\\
\texttt{$\rightarrow$ remediation\_explainer $\rightarrow$ content\_guardrail $\rightarrow$ report\_generator $\rightarrow$ notifier $\rightarrow$ END}
\end{center}

La figure \ref{fig:architecture_ccu} illustre ce pipeline agentique.

\begin{figure}[H]
    \centering
    \begin{tikzpicture}[
    node distance=0.9cm and 0.6cm,
    box/.style={rectangle, rounded corners, draw=iNavy, fill=fondBleu, text=iNavy, font=\small\bfseries, minimum width=2.8cm, minimum height=0.9cm, align=center},
    smallbox/.style={rectangle, rounded corners, draw=RoyalBlue, fill=white, text=black, font=\scriptsize, minimum width=2.6cm, minimum height=0.7cm, align=center},
    arrow/.style={-{Stealth[length=2mm]}, thick, iNavy},
    label/.style={font=\scriptsize\itshape, text=gray}
]

% Main flow nodes
\node[box] (start) {START};
\node[box, below=of start] (intake) {intake\_parser};
\node[box, below=2.2cm of intake] (root) {root\_cause\_reasoner};
\node[box, below=of root] (ticket) {ticket\_manager};
\node[box, below=of ticket] (remediation) {remediation\_explainer};
\node[box, below=of remediation] (guardrail) {content\_guardrail};
\node[box, below=of guardrail] (report) {report\_generator};
\node[box, below=of report] (notifier) {notifier};
\node[box, below=of notifier] (end) {END};

% Collector nodes (parallel)
\node[smallbox, left=1.2cm of root, yshift=1.6cm] (logs) {logs\_investigator};
\node[smallbox, left=of root, yshift=1.6cm] (context) {context\_agent};
\node[smallbox, right=of root, yshift=1.6cm] (rag) {rag\_ticket\_search};

% Arrows main flow
\draw[arrow] (start) -- (intake);
\draw[arrow] (intake) -- ++(0,-0.8) -| (logs);
\draw[arrow] (intake) -- ++(0,-0.8) -| (context);
\draw[arrow] (intake) -- ++(0,-0.8) -| (rag);
\draw[arrow] (logs) |- (root);
\draw[arrow] (context) |- (root);
\draw[arrow] (rag) |- (root);
\draw[arrow] (root) -- (ticket);
\draw[arrow] (ticket) -- (remediation);
\draw[arrow] (remediation) -- (guardrail);
\draw[arrow] (guardrail) -- (report);
\draw[arrow] (report) -- (notifier);
\draw[arrow] (notifier) -- (end);

% Collector label
\node[label, below=0.15cm of intake, xshift=0cm] {agents collecteurs en parallèle};

% Side outputs
\node[right=0.6cm of report, font=\scriptsize, align=left] {PDF/HTML};
\node[right=0.6cm of notifier, font=\scriptsize, align=left] {email +\\note Zammad};

\end{tikzpicture}

    \caption{Pipeline agentique du diagnostic technique CCU}
    \label{fig:architecture_ccu}
\end{figure}

\section{Modèle de données partagé}
L'état du graphe est matérialisé par la classe \texttt{GraphState} définie dans \texttt{shared/state.py}. Elle utilise \textbf{Pydantic} pour valider les données à chaque transition. Les champs principaux sont :

\begin{itemize}
    \setlength\itemsep{0.1em}
    \setlength\parskip{0pt}
    \setlength\parsep{0pt}
    \item \texttt{incident} : description brute de l'incident.
    \item \texttt{parsed\_incident} : structure normalisée (client, service, commande, type, priorité).
    \item \texttt{logs}, \texttt{customer\_context}, \texttt{similar\_tickets} : résultats des agents collecteurs.
    \item \texttt{root\_cause} : cause racine et niveau de confiance.
    \item \texttt{remediation\_explanation} : explication de la remédiation proposée.
    \item \texttt{ticket\_mapping} : mapping ou création du ticket.
    \item \texttt{sanitized\_*} : champs filtrés par le guardrail.
    \item \texttt{report\_path}, \texttt{email\_sent}, \texttt{zammad\_note\_added} : sorties finales.
    \item \texttt{traces} : journal d'audit de l'exécution.
\end{itemize}

La classe \texttt{ParsedIncident} formalise les attributs extraits lors de l'intake :

\begin{table}[H]
    \centering
    \renewcommand{\arraystretch}{1.2}
    \begin{tabular}{|l|p{10cm}|}
        \hline
        \rowcolor{iNavy} \textbf{\color{white}Attribut} & \textbf{\color{white}Description} \\ \hline
        service\_id & Identifiant du service concerné (ex. svc-fiber-12345) \\ \hline
        order\_id & Identifiant de la commande (ex. ord-2026-001) \\ \hline
        customer\_id & Identifiant du client (ex. acc-12345) \\ \hline
        incident\_type & Type d'incident : fibre, sim, mobile, lan, billing, autre \\ \hline
        description & Texte décrivant l'incident \\ \hline
        priority & Priorité P1, P2 ou P3 \\ \hline
    \end{tabular}
    \caption{Attributs de l'incident normalisé}
    \label{tab:parsed_incident}
\end{table}

\section{Les agents du pipeline}

\subsection{Intake parser}
L'agent \texttt{intake\_parser} a pour mission de normaliser l'incident brut. Il combine une phase de parsing heuristique par expressions régulières et une validation/ajustement par LLM structuré lorsque le modèle est disponible. En mode \texttt{MOCK\_LLM=true}, le système fonctionne de manière déterministe sans dépendance externe.

\subsection{Collecteurs parallèles}
Les trois agents collecteurs s'exécutent en parallèle dans le nœud \texttt{collectors} :

\begin{itemize}
    \setlength\itemsep{0.1em}
    \setlength\parskip{0pt}
    \setlength\parsep{0pt}
    \item \texttt{logs\_investigator} : synthétise les logs réseau et applicatifs pertinents.
    \item \texttt{context\_agent} : récupère le contexte client et contractuel dans le CRM.
    \item \texttt{rag\_ticket\_search} : recherche des tickets historiques similaires dans Neo4j.
\end{itemize}

Le parallélisme est mis en œuvre via \texttt{asyncio.gather}, avec un fallback \texttt{ThreadPoolExecutor} en cas de boucle déjà active.

\subsection{Root cause reasoner}
L'agent \texttt{root\_cause\_reasoner} corrèle les sources collectées pour établir une cause racine. Il retourne un niveau de confiance (faible, moyenne, forte) et une liste de \texttt{source\_ids}. Si aucune source suffisante n'est disponible, il retourne une cause \texttt{indéterminée} avec une confiance faible, ce qui constitue un mécanisme de refus d'hallucination.

\subsection{Ticket manager}
L'agent \texttt{ticket\_manager} détermine s'il faut créer un nouveau ticket ou rattacher l'incident à un ticket existant. Il s'appuie sur le client, le service et les tickets similaires déjà présents.

\subsection{Remediation explainer}
L'agent \texttt{remediation\_explainer} produit un texte informatif décrivant les actions recommandées. Conformément au principe de sécurité du projet, ces actions ne sont jamais exécutées automatiquement.

\subsection{Content guardrail}
Le guardrail \texttt{pii\_sanitizer} filtre les informations personnelles identifiables (PII) avant toute génération de rapport ou notification externe. Il classe également les actions selon une liste blanche pour refuser tout comportement non autorisé.

\subsection{Report generator et notifier}
L'agent \texttt{report\_generator} produit un rapport PDF/HTML à partir d'un template Jinja2, avec WeasyPrint ou un fallback HTML. L'agent \texttt{notifier} envoie l'email et ajoute une note interne dans Zammad.

\section{Graphe de connaissances Neo4j}

\subsection{Modèle de données graphe}
Le graphe Neo4j stocke les entités métiers et leurs relations. Il sert à la fois de base de connaissances structurée et de support à la recherche vectorielle. Les principaux nœuds sont : \texttt{Client}, \texttt{Subscription}, \texttt{Order}, \texttt{Product}, \texttt{Ticket} et \texttt{LogEvent}.

La figure \ref{fig:data_model_ccu} représente le modèle de données graphe.

\begin{figure}[H]
    \centering
    \begin{tikzpicture}[
    node distance=1.6cm,
    entity/.style={rectangle, rounded corners, draw=iNavy, fill=fondBleu, text=iNavy, font=\small\bfseries, minimum width=2.4cm, minimum height=0.8cm, align=center},
    relation/.style={-{Stealth[length=2mm]}, thick, iNavy},
    relabel/.style={font=\scriptsize, text=gray, fill=white, inner sep=1pt}
]

% Nodes
\node[entity] (client) {Client};
\node[entity, right=of client] (sub) {Subscription};
\node[entity, right=of sub] (order) {Order};
\node[entity, below=of client] (product) {Product};
\node[entity, below=of sub] (ticket) {Ticket};
\node[entity, below=of order] (log) {LogEvent};

% Relationships
\draw[relation] (client) -- node[relabel, above] {HAS\_SUBSCRIPTION} (sub);
\draw[relation] (sub) -- node[relabel, above] {SUBSCRIBED\_TO} (product);
\draw[relation] (client) -- node[relabel, left] {PLACED} (order);
\draw[relation] (order) -- node[relabel, right] {FOR\_PRODUCT} (product);
\draw[relation] (sub) -- node[relabel, left] {HAS\_TICKET} (ticket);
\draw[relation] (order) -- node[relabel, right] {GENERATED} (log);
\draw[relation, dashed] (ticket) -- node[relabel, below] {RELATED\_TO\_PRODUCT} (product);
\draw[relation, dashed, bend left=30] (ticket) -- node[relabel, above] {SIMILAR\_TO} (ticket);

\end{tikzpicture}

    \caption{Modèle de données du graphe Neo4j}
    \label{fig:data_model_ccu}
\end{figure}

\subsection{Contraintes et index vectoriel}
Le fichier \texttt{graph/schema.cypher} définit les contraintes d'unicité sur chaque entité et crée un index vectoriel sur les embeddings des tickets. Cet index permet d'effectuer une recherche de similarité cosine nativement dans Neo4j.

\begin{lstlisting}[language=SQL, caption={Extrait du schema.cypher}, label={lst:schema_cypher}]
CREATE CONSTRAINT ticket_id IF NOT EXISTS
FOR (t:Ticket) REQUIRE t.id IS UNIQUE;

CREATE VECTOR INDEX ticket_embeddings IF NOT EXISTS
FOR (t:Ticket) ON (t.embedding)
OPTIONS {indexConfig: {
    `vector.dimensions`: __VECTOR_DIM__,
    `vector.similarity_function`: 'cosine'
}};
\end{lstlisting}

\subsection{Recherche GraphRAG}
La requête principale \texttt{search\_similar\_incidents}, définie dans \texttt{graph/queries.py}, combine deux mécanismes :

\begin{enumerate}
    \item \textbf{Recherche vectorielle} : l'incident courant est vectorisé par un provider d'embedding, puis Neo4j retourne les tickets les plus proches en termes de cosine similarity.
    \item \textbf{Filtre de proximité de graphe} : seuls les tickets liés au même produit ou au même type de commande que l'incident courant sont conservés.
\end{enumerate}

Cette approche \textbf{GraphRAG} améliore la précision par rapport à une simple recherche sémantique en intégrant la structure relationnelle du graphe.

\subsection{Provider d'embeddings}
L'interface \texttt{EmbeddingProvider} définie dans \texttt{graph/embedding\_provider.py} permet de changer de provider d'embedding (Ollama, sentence-transformers, OpenAI, Voyage AI) sans modifier le reste du code. Le provider Ollama est utilisé par défaut avec le modèle \texttt{mxbai-embed-large:latest}.

\section{API REST et interface utilisateur}

\subsection{FastAPI et streaming SSE}
L'API est développée avec \textbf{FastAPI}. La route \texttt{/diagnose/stream} retourne un flux d'événements Server-Sent Events (SSE), chaque événement correspondant à un nœud traversé par le graphe LangGraph. Cette approche permet au client (chatbot ou outil externe) de suivre la progression du diagnostic en temps réel.

\subsection{Chatbot Next.js}
L'interface utilisateur est un chatbot développé avec \textbf{Next.js}. Il affiche la trace d'exécution des agents, le badge de création du ticket et la carte de rapport avec le lien de téléchargement PDF.

\section*{Conclusion}
Ce chapitre a présenté l'architecture agentique du projet : un pipeline LangGraph modulaire, un état partagé validé par Pydantic, des agents spécialisés et un graphe de connaissances Neo4j exploitant la recherche GraphRAG. Le chapitre suivant détaille l'implémentation technique de ces composants.
